\section{Motor de Inteligencia Artificial}

\subsection{Gestión del historial y prevención del desbordamiento de contexto}

Los modelos de lenguaje tienen un límite en la cantidad de texto que pueden 
recordar a la vez (lo que se conoce como ventana de contexto, que en nuestro 
caso es de 8192 tokens). Como las instrucciones iniciales y la explicación de 
las herramientas ya ocupan bastante espacio, es necesario controlar cuánto 
historial de la conversación le enviamos para no saturarlo.

Para solucionar esto, la clase \texttt{Agente} (en \texttt{agent.py}) incluye una 
función llamada \texttt{recortar\_historial} que limita la memoria a los últimos 
24 mensajes. Este recorte no se hace de cualquier manera, sino que el sistema siempre conserva el primer mensaje (el 
\textit{System Prompt}) y busca el inicio de una intervención del usuario 
(\texttt{role: "user"}) para cortar a partir de ahí.   

De esta forma, nos aseguramos de no partir por la mitad una interacción entre el 
modelo y una herramienta. Si cortáramos mal, el modelo podría recibir la respuesta
de una herramienta sin saber por qué la había pedido, lo que provocaría un error 
al intentar generar el siguiente mensaje

\subsection{System Prompt dinámico y reglas de comportamiento}

El comportamiento del asistente y las reglas que debe seguir se definen en un 
\textit{System Prompt} que construimos en el archivo \texttt{config.py}.

Este texto no es estático, sino que se genera cada vez que arranca la aplicación 
para poder incluir información actualizada. Por ejemplo, le pasamos la fecha de 
hoy usando \texttt{date.today().isoformat()} para que el modelo entienda a qué 
nos referimos cuando le preguntamos por los gastos de este mes o 
la semana pasada. También le adjuntamos la estructura de la base de datos 
(\texttt{ESQUEMA\_BD}) para que sepa qué tablas y columnas existen cuando tenga 
que crear consultas SQL.



Además, en este \textit{prompt} le damos instrucciones claras sobre cómo debe 
expresarse. Le pedimos que trate al usuario de tú, que formatee el dinero al 
estilo español (por ejemplo, "1.234,56 €") y le prohibimos usar formato 
Markdown como negritas o cursivas. Esto último es muy importante para que el 
sistema que lee el texto en voz alta lo haga de forma natural, sin tropezar 
leyendo símbolos.   

\subsection{Uso de Tool Calling }

Para que el modelo pueda interactuar con la aplicación, utilizamos la 
funcionalidad de \textit{Tool Calling}. En el archivo \texttt{tools.py} 
definimos en formato JSON una lista con todo lo que el asistente puede invocar: 
consultar datos, preparar un envío de Bizum\dots

Para ejecutar el modelo de lenguaje en local, garantizar la privacidad de los datos
y no depender de suscripciones o servicios en la nube hemos utilizado Ollama. 

El problema que tenemos es que las herramientas en \texttt{tools.py} las están 
definidas usando la estructura de Anthropic, pero Ollama no soporta directamente 
este formato (esto pues empezamos el proyecto pensando que usariamos API de Claude 
pero finalmente decidimos priorizar el uso de un modelo local). 

Para solucionarlo, en \texttt{agent.py} está la función 
\texttt{adaptar\_herramientas\_a\_openai}, que hace de puente y traduce 
automáticamente nuestras herramientas al formato correcto antes de enviárselas al 
modelo, permitiendo que la conexión funcione sin problemas.

El proceso completo de interacción funciona como se ilustra en la
Figura~\ref{fig:flujo_agente}:
el usuario hace una pregunta, y el modelo evalúa si necesita usar alguna herramienta para responder. Si es así, interrumpe su respuesta y devuelve un mensaje especial pidiendo ejecutarla. Nuestro \emph{backend} recibe esa petición, ejecuta el código en Python (por ejemplo, hace la consulta a la base de datos) y le devuelve el resultado al modelo añadiendo un mensaje con el rol \texttt{tool}. Finalmente, el modelo lee ese resultado y ya dispone de los datos concretos para redactar la respuesta que escuchará el usuario.

\begin{figure}[htbp]
\centering
\begin{tikzpicture}[
    font=\small,
    caja/.style={draw=#1, fill=#1!8, rounded corners=3pt, align=center,
                 minimum width=3.4cm, minimum height=1.15cm, inner sep=3pt},
    caja/.default=subsectionblue,
    hist/.style={draw=gray!60, fill=gray!8, rounded corners=2pt, align=left,
                 font=\scriptsize\ttfamily, inner sep=4pt, minimum width=4.6cm},
    flecha/.style={-{Stealth[length=2mm]}, thick, draw=gray!70!black},
    et/.style={font=\scriptsize, text=gray!55!black, align=center, inner sep=2pt},
]

% --- Actores ---------------------------------------------------------------
\node[caja]        (ag)  at (0,0)    {\texttt{agent.py}\\[-1pt]\scriptsize bucle agéntico};
\node[caja=ugrred] (llm) at (7.6,0)  {Modelo\\[-1pt]\scriptsize \texttt{qwen3:8b}};
\node[caja=gray]   (fn)  at (0,-3.8) {Función Python\\[-1pt]\scriptsize \texttt{ejecutar\_tool}};

% --- Ciclo -----------------------------------------------------------------
\draw[flecha] ([yshift=3mm]ag.east) -- node[et, above]
      {1. mensajes + esquemas\\de las 6 herramientas} ([yshift=3mm]llm.west);
\draw[flecha] ([yshift=-3mm]llm.west) -- node[et, below]
      {2. \texttt{tool\_calls}: nombre\\y argumentos en JSON} ([yshift=-3mm]ag.east);
\draw[flecha] (ag) -- node[et, left] {3. ejecuta} (fn);
\draw[flecha] ([xshift=6mm]fn.north) -- node[et, right]
      {4. resultado como\\mensaje de rol \texttt{tool}} ([xshift=6mm]ag.south);

\draw[flecha, draw=subsectionblue] (llm.north) .. controls (7.6,2.0) and (0,2.0) ..
      node[et, above, text=subsectionblue, pos=0.5] {se repite hasta que el modelo responde sin pedir herramientas (máx.\ 8 vueltas)}
      (ag.north);

% --- Historial resultante --------------------------------------------------
\node[hist, anchor=north west] (h) at (4.0,-2.6)
      {\textcolor{gray}{system}: reglas y esquema\\
       \textcolor{gray}{user}: ¿cuánto gasté en gasolina?\\
       \textcolor{gray}{assistant}: \textit{tool\_calls}\\
       \textcolor{gray}{tool}: \{"filas": [[84.2]]\}\\
       \textcolor{gray}{assistant}: Has gastado 84,20 €};
\node[et, anchor=south west] at (4.0,-2.5) {El historial que se reenvía en cada vuelta};

\end{tikzpicture}
\caption{Ciclo de \textit{tool calling}. El modelo nunca ejecuta nada: devuelve
el nombre de la herramienta y sus argumentos, el \emph{backend} ejecuta la
función en Python y le devuelve el resultado como un mensaje más del historial.
Con ese dato ya concreto, el modelo redacta la respuesta final.}
\label{fig:flujo_agente}
\end{figure}


